% Generated draft: requires mathematical and source review.
\documentclass[11pt]{article}
\usepackage[margin=25mm]{geometry}
\usepackage{fontspec}
\setmainfont{Latin Modern Roman}
\usepackage{amsmath,amssymb,mathtools,unicode-math}
\setmathfont{Latin Modern Math}
\usepackage{enumitem,needspace,xurl,hyperref}
\setlength{\parindent}{0pt}
\setlength{\parskip}{5pt}
\setlength{\emergencystretch}{3em}
\newcommand{\N}{\mathbb{N}}
\newcommand{\Z}{\mathbb{Z}}
\newcommand{\Q}{\mathbb{Q}}
\newcommand{\R}{\mathbb{R}}
\newcommand{\C}{\mathbb{C}}
\newcommand{\F}{\mathbb{F}}
\newcommand{\A}{\mathbb{A}}
\newcommand{\PP}{\mathbb{P}}
\newcommand{\E}{\mathbb{E}}
\newcommand{\eps}{\varepsilon}
\newcommand{\dd}{\,\mathrm{d}}
\begin{document}
\section*{UnsolvedMath ID: 800672}
\textbf{Source problem code:} catalogue entry
\subsection{Definitions and mathematical statement}
For a topological space \(Y\) and \(y\in Y\), a family \(\mathcal P\subseteq\mathcal P(Y)\) is a \(\pi\)-network at \(y\) if every neighborhood \(V\) of \(y\) contains some \(P\in\mathcal P\). The space \(Y\) has the Pytkeev property if, for every \(y\in Y\) and every \(A\subseteq Y\) with \(y\in\overline A\setminus A\), there is a countable \(\pi\)-network \(\mathcal P\) at \(y\) such that every \(P\in\mathcal P\) is an infinite subset of \(A\).

Under the usual interpretation, \(C_p(X)\) denotes the space of continuous functions \(X\to\R\), with the topology inherited from the product \(\R^X\). The supplied record does not specify the class of topological spaces \(X\) for which the cited characterization is intended. This interpretation of \(C_p(X)\) does not resolve that missing hypothesis.

Using the standard convention, an \(\omega\)-cover of \(X\) is a family \(\mathcal U\) of subsets of \(X\) such that \(X\notin\mathcal U\) and, for every finite \(F\subseteq X\), there is \(U\in\mathcal U\) with \(F\subseteq U\). It is open if every member is open in \(X\). The supplied record does not define the additional condition called \(\omega\)-shrinkability; its intended definition must be checked in the cited source.

The record attributes to [SakaiRez] a characterization of the Pytkeev property of \(C_p(X)\) by requiring every \(\omega\)-shrinkable open \(\omega\)-cover \(\mathcal U\) to satisfy the following displayed condition, reproduced literally:
\[
\exists (\mathcal U_n)_{n\in\N}\quad
\left[
 (\forall n\in\N)\;\mathcal U_n\subseteq\mathcal U
 \quad\text{and}\quad
 \bigcap_{n\in\N}\mathcal U_n
 \text{ is an }\omega\text{-cover of }X
\right]. \tag{*}
\]
The source question, attributed to [SakaiRez, Question 2.8], asks whether \(\omega\)-shrinkability can be omitted from that characterization: for every space \(X\) in its intended domain, is \(C_p(X)\) a Pytkeev space if and only if every open \(\omega\)-cover \(\mathcal U\) satisfies the intended cover condition?

There is an unresolved formulation issue. In \((*)\), the intersection is an intersection of families, and no infinitude requirement is imposed on the subfamilies \(\mathcal U_n\). Literally, \((*)\) holds for every \(\omega\)-cover by taking \(\mathcal U_n=\mathcal U\) for all \(n\). Thus the supplied formula does not specify a substantive additional cover property. Whether a family of intersections of cover members was intended, and whether the subfamilies were required to be infinite, cannot be settled from this record. No corrected formula or missing hypothesis is asserted here.
\subsection{Short English statement}
The question asks whether the cover characterization of the Pytkeev property of \(C_p(X)\), attributed to [SakaiRez], remains valid for all open \(\omega\)-covers after removing \(\omega\)-shrinkability. The supplied cover condition is automatic as written, so its intended formula, the definition of shrinkability, and the hypotheses on \(X\) require clarification.
\subsection{Sources}
\begin{itemize}
\item Local UnsolvedMath registry, permanent ID \(800672\), exact source code \texttt{catalogue entry}; title: Selection Principles Milestone Problems --- Problem 45 ([SakaiRez, Question 2.8]). The registry lists Tsaban as proposer, the year 2003, and status \texttt{open}. This status is supplied metadata, not an independently verified mathematical or literature assessment.
\item Supplied source: Tsaban, \emph{Selection Principles in Mathematics: A Milestone of Open Problems} (2003), Problem 45, \url{https://arxiv.org/abs/math/0312182}. The record reports access on 2026-10-16 and extraction by \texttt{arxiv\_latex}; these are supplied provenance details, not a claim of independent access or verification.
\item The record attributes the characterization to [SakaiRez] and the question to [SakaiRez, Question 2.8], also describing the characterization as Sakai's. No full bibliographic entry or direct link for [SakaiRez] is supplied, so that underlying reference has not been identified from the provided evidence.
\end{itemize}
\end{document}
